%%%%%%%%%%%%%%%%%%%%%%%%%%%%%%%%%%%%%%%%%%%%%%%%%%%%%%%%%%%%%%%%%%%%%%
% Lớp 8 | Chương III. Tứ giác | Tứ giác
% Dạng 8H3?10-2: Tính góc dựa vào tổng các góc của tứ giác
%%%%%%%%%%%%%%%%%%%%%%%%%%%%%%%%%%%%%%%%%%%%%%%%%%%%%%%%%%%%%%%%%%%%%%


% nguồn: De_AI_Toan8_cuoi_hoc_ki_i_1 (nhập 21/11/2026)
\begin{bt}%[8H3N10-2]%[ma:a194fa1c]
	Cho tứ giác $ABCD$ có $\widehat{A} = 85^\circ$, $\widehat{B} = 115^\circ$, $\widehat{C} = 70^\circ$. Tính số đo góc $D$.
	\loigiai{
		Xét tứ giác $ABCD$, theo định lí tổng các góc của một tứ giác, ta có: \\
$\widehat{A} + \widehat{B} + \widehat{C} + \widehat{D} = 360^\circ$. \\
Thay số: $85^\circ + 115^\circ + 70^\circ + \widehat{D} = 360^\circ$. \\
$270^\circ + \widehat{D} = 360^\circ \Rightarrow \widehat{D} = 360^\circ - 270^\circ = 90^\circ$. \\
Vậy $\widehat{D} = 90^\circ$.
	}
\end{bt}

% nguồn: De_AI_Toan8_cuoi_hoc_ki_i_1 (nhập 21/11/2026)
\begin{bt}%[8H3N10-2]%[ma:b9ece1db]
	Cho tứ giác $ABCD$ có $\widehat{A} = 100^\circ$, $\widehat{B} = 90^\circ$, $\widehat{C} = 80^\circ$.
a) Tính số đo góc $D$.
b) Tính số đo góc ngoài tại đỉnh $A$ và đỉnh $D$ của tứ giác $ABCD$.
	\loigiai{
		a) Trong tứ giác $ABCD$, ta có $\widehat{A} + \widehat{B} + \widehat{C} + \widehat{D} = 360^\circ$. \\
Suy ra $\widehat{D} = 360^\circ - (100^\circ + 90^\circ + 80^\circ) = 360^\circ - 270^\circ = 90^\circ$. \\
b) Vì góc ngoài kề bù với góc trong tại cùng một đỉnh nên: \\
Góc ngoài tại đỉnh $A$ bằng $180^\circ - \widehat{A} = 180^\circ - 100^\circ = 80^\circ$. \\
Góc ngoài tại đỉnh $D$ bằng $180^\circ - \widehat{D} = 180^\circ - 90^\circ = 90^\circ$.
	}
\end{bt}

% nguồn: De_AI_Toan8_cuoi_hoc_ki_i_1 (nhập 21/11/2026)
\begin{bt}%[8H3N10-2]%[ma:d2cefabf]
	Cho tứ giác $MNPQ$ có $\widehat{M} = 110^\circ$, $\widehat{N} = 75^\circ$ và $\widehat{P} = \widehat{Q}$. Tính số đo góc $P$ và góc $Q$.
	\loigiai{
		Trong tứ giác $MNPQ$, ta có $\widehat{M} + \widehat{N} + \widehat{P} + \widehat{Q} = 360^\circ$. \\
Thay $\widehat{M} = 110^\circ$, $\widehat{N} = 75^\circ$ và $\widehat{Q} = \widehat{P}$, ta được: \\
$110^\circ + 75^\circ + 2\widehat{P} = 360^\circ \Rightarrow 2\widehat{P} = 360^\circ - 185^\circ = 175^\circ$. \\
Suy ra $\widehat{P} = \widehat{Q} = \dfrac{175^\circ}{2} = 87{,}5^\circ$.
	}
\end{bt}

% nguồn: De_AI_Toan8_cuoi_hoc_ki_i_1 (nhập 21/11/2026)
\begin{bt}%[8H3N10-2]%[ma:913952d7]
	Cho tứ giác $ABCD$ có số đo các góc $\widehat{A}$, $\widehat{B}$, $\widehat{C}$, $\widehat{D}$ tỉ lệ thuận với $1$; $2$; $3$; $4$. Tính số đo mỗi góc của tứ giác $ABCD$.
	\loigiai{
		Theo đề bài ta có $\dfrac{\widehat{A}}{1} = \dfrac{\widehat{B}}{2} = \dfrac{\widehat{C}}{3} = \dfrac{\widehat{D}}{4}$ và $\widehat{A} + \widehat{B} + \widehat{C} + \widehat{D} = 360^\circ$. \\
Áp dụng tính chất dãy tỉ số bằng nhau: \\
$\dfrac{\widehat{A}}{1} = \dfrac{\widehat{B}}{2} = \dfrac{\widehat{C}}{3} = \dfrac{\widehat{D}}{4} = \dfrac{\widehat{A} + \widehat{B} + \widehat{C} + \widehat{D}}{1 + 2 + 3 + 4} = \dfrac{360^\circ}{10} = 36^\circ$. \\
Suy ra $\widehat{A} = 36^\circ$, $\widehat{B} = 72^\circ$, $\widehat{C} = 108^\circ$, $\widehat{D} = 144^\circ$.
	}
\end{bt}

% nguồn: De_AI_Toan8_cuoi_hoc_ki_i_1 (nhập 21/11/2026)
\begin{bt}%[8H3N10-2]%[ma:9cf53df8]
	Cho tứ giác $ABCD$ có $\widehat{A} = 60^\circ$, $\widehat{B} = 110^\circ$ và $\widehat{C} - \widehat{D} = 30^\circ$. Tính số đo góc $C$ và góc $D$.
	\loigiai{
		Trong tứ giác $ABCD$, ta có $\widehat{A} + \widehat{B} + \widehat{C} + \widehat{D} = 360^\circ$. \\
Suy ra $\widehat{C} + \widehat{D} = 360^\circ - (\widehat{A} + \widehat{B}) = 360^\circ - (60^\circ + 110^\circ) = 190^\circ$. \\
Kết hợp với $\widehat{C} - \widehat{D} = 30^\circ$, ta được: \\
$\widehat{C} = \dfrac{190^\circ + 30^\circ}{2} = 110^\circ$ và $\widehat{D} = 110^\circ - 30^\circ = 80^\circ$.
	}
\end{bt}

% nguồn: De_AI_Toan8_cuoi_hoc_ki_i_1 (nhập 21/11/2026)
\begin{bt}%[8H3N10-2]%[ma:32d2db34]
	Cho tứ giác $ABCD$ có $\widehat{A} = x$, $\widehat{B} = x + 20^\circ$, $\widehat{C} = x + 40^\circ$, $\widehat{D} = x + 60^\circ$. Tìm $x$ và tính số đo các góc của tứ giác $ABCD$.
	\loigiai{
		Vì tổng bốn góc trong tứ giác $ABCD$ bằng $360^\circ$ nên: \\
$x + (x + 20^\circ) + (x + 40^\circ) + (x + 60^\circ) = 360^\circ$. \\
$4x + 120^\circ = 360^\circ \Rightarrow 4x = 240^\circ \Rightarrow x = 60^\circ$. \\
Vậy $\widehat{A} = 60^\circ$, $\widehat{B} = 80^\circ$, $\widehat{C} = 100^\circ$, $\widehat{D} = 120^\circ$.
	}
\end{bt}

% nguồn: De_AI_Toan8_cuoi_hoc_ki_i_1 (nhập 21/11/2026)
\begin{bt}%[8H3N10-2]%[ma:3ddf58cf]
	Cho tứ giác $ABCD$ có các góc ngoài tại đỉnh $A$, $B$, $C$ lần lượt là $70^\circ$, $85^\circ$, $95^\circ$.
a) Tính số đo các góc trong $\widehat{A}$, $\widehat{B}$, $\widehat{C}$.
b) Tính số đo góc trong và góc ngoài tại đỉnh $D$.
	\loigiai{
		a) Vì góc trong và góc ngoài tại mỗi đỉnh là hai góc kề bù nên: \\
$\widehat{A} = 180^\circ - 70^\circ = 110^\circ$; \\
$\widehat{B} = 180^\circ - 85^\circ = 95^\circ$; \\
$\widehat{C} = 180^\circ - 95^\circ = 85^\circ$. \\
b) Trong tứ giác $ABCD$, ta có $\widehat{D} = 360^\circ - (110^\circ + 95^\circ + 85^\circ) = 360^\circ - 290^\circ = 70^\circ$. \\
Góc ngoài tại đỉnh $D$ bằng $180^\circ - 70^\circ = 110^\circ$.
	}
\end{bt}

% nguồn: De_AI_Toan8_cuoi_hoc_ki_i_1 (nhập 21/11/2026)
\begin{bt}%[8H3N10-2]%[ma:bc73212d]
	Cho tứ giác $ABCD$ có $\widehat{A} = 90^\circ$, $\widehat{B} = 120^\circ$, $\widehat{C} = 80^\circ$. Các tia phân giác của góc $C$ và góc $D$ cắt nhau tại $I$.
a) Tính số đo góc $D$.
b) Tính số đo góc $CID$.
	\loigiai{
		a) Trong tứ giác $ABCD$, ta có $\widehat{D} = 360^\circ - (\widehat{A} + \widehat{B} + \widehat{C}) = 360^\circ - (90^\circ + 120^\circ + 80^\circ) = 70^\circ$. \\
b) Vì $CI$, $DI$ lần lượt là tia phân giác của $\widehat{C}$ và $\widehat{D}$ nên: \\
$\widehat{ICD} = \dfrac{\widehat{C}}{2} = \dfrac{80^\circ}{2} = 40^\circ$ và $\widehat{IDC} = \dfrac{\widehat{D}}{2} = \dfrac{70^\circ}{2} = 35^\circ$. \\
Xét tam giác $CID$, ta có $\widehat{CID} = 180^\circ - (\widehat{ICD} + \widehat{IDC}) = 180^\circ - (40^\circ + 35^\circ) = 105^\circ$.
	}
\end{bt}

% nguồn: chuyen_dang (nhập 22/11/2026)
\begin{bt}%[8H3H10-2]%[ma:28668c45]
    Cho tam giác $ABC$ có $\widehat{A}=100^{\circ}$, $\widehat{C}=35^{\circ}$ và $b=10$. Tính $a$, $c$, $R$ và số đo góc $B$.
    \loigiai{
        Ta có $\widehat{B}=180^{\circ}-\left(\widehat{A}+\widehat{C}\right)=180^{\circ}-\left(100^{\circ}+35^{\circ}\right)=45^{\circ}$.\\
        Áp dụng Định lí sin, ta có $\dfrac{a}{\sin 100^{\circ}}=\dfrac{10}{\sin 45^{\circ}}=\dfrac{c}{\sin 35^{\circ}}=2R$.\\
        Suy ra $\heva{&a=\dfrac{10}{\sin 45^{\circ}}\cdot\sin 100^{\circ}=10\sqrt{2}\cdot\sin 100^{\circ}\approx 13{,}93\\&c=\dfrac{10}{\sin 45^{\circ}}\cdot\sin 35^{\circ}=10\sqrt{2}\cdot\sin 35^{\circ}\approx 8{,}11\\&R=\dfrac{10}{2\sin 45^{\circ}}=5\sqrt{2}.}$
    }
\end{bt}

% nguồn: chuyen_dang (nhập 22/11/2026)
\begin{bt}%[8H3V10-2]%[ma:86ec9cc0]
	Cho hình chóp $S. A B C D$ có đáy là hình bình hành tâm $O$. Biết $\widehat{S A D}=\widehat{S C  D}=90^{\circ}$. Tìm số đo của góc $(S B, A C)$
	\loigiai{ 
	\immini{ 
		Gọi  $M$ là trung điểm cạnh $S D$.\\
		Khi đó $O M$ là đường trung bình của tam giác $S B D$ nên $O M \parallel  S B$.\\
		Do vậy $(S  B, A C)=(O M, A C)$.\\
		Tam giác $S A  D$ vuông tại $A$ có trung tuyến $A M=\dfrac{S D}{2}$;\\
		Tam giác $S C D$ vuông tại $C$  có trung tuyến $C M=\dfrac{S D}{2}$.\\
		Suy ra tam giác $M  A C$ cân tại $M$, mà $O$ là trung điểm $A C$ nên $O M \perp A C$ hay $(O M, A C)=90^{\circ}$. Do đó $(S B, A C)=90^{\circ}$ hay $S B \perp A C$.
		} 
		{ 
		\begin{tikzpicture}[scale=0.9]
				 \tkzDefPoints{0/0/D,-2/-2/C,2.5/-2/B}
				\coordinate (A)  at ($(B)+(D)-(C)$);
				\draw (C)--(B)--(A); 
				 \draw[dashed] (C)--(D)--(A);
				\coordinate (O) at ($(A)!0.5!(C)$); 
				\coordinate (S)  at ($(O)+(-0.5,5)$);
				\coordinate  (M) at ($(S)!0.5!(A)$);
				\draw[dashed] (O)--(S)--(D)--(B) 
					 (A)--(C)
					(D)--(M)--(O); 
				\draw (S)--(A) 
					(S)--(B)--(M) 
					 (S)--(C);
				\tkzMarkRightAngle  (A,O,D)
				 \tkzMarkRightAngle (S,D,A)
				\tkzMarkRightAngle (S,B,A) 
				\begin{footnotesize}
					\draw (S) node[above]{$S$}; 
					 \draw (C) node[left]{$B$};
					 \draw (D) node[left]{$A$};
					\draw (A) node[right]{$D$}; 
					\draw  (B) node[right]{$C$};
					\draw  (M) node[right]{$M$};
					\draw  (O) node[below]{$O$};
				\end{footnotesize}
			\end{tikzpicture}
		} 
	} 
\end{bt}

% nguồn: chuyen_dang (nhập 22/11/2026)
\begin{ex}%[8H3V10-2]%[ma:1b2f549f]
	Cho hình chóp $S. A B C D$ có đáy là hình bình hành tâm $O$. Biết $\widehat{S A D}=\widehat{S C D}=90^{\circ}$. Tìm số đo  của góc $(S B, A C)$
	 \loigiai{
	\immini{ 
		Gọi $M$  là trung điểm cạnh $S D$.\\
		Khi đó $O M$ là đường trung bình của tam giác $S B D$ nên  $O M \parallel S B$.\\
		 Do vậy $(S B, A C)=(O M, A C)$.\\
		Tam giác $S A D$ vuông tại $A$  có trung tuyến $A M=\dfrac{S D}{2}$;\\
		Tam giác $S C D$ vuông tại  $C$ có trung tuyến $C M=\dfrac{S D}{2}$.\\
		Suy ra tam  giác $M A C$ cân tại $M$, mà $O$ là trung điểm $A C$ nên $O M \perp A C$ hay $(O M, A C)=90^{\circ}$. Do đó $(S B, A C)=90^{\circ}$ hay $S B \perp A C$.
		} 
		 {
		\begin{tikzpicture}[scale=0.9]
				 \tkzDefPoints{0/0/D,-2/-2/C,2.5/-2/B}
				\coordinate (A) at ($(B)+(D)-(C)$); 
				 \draw (C)--(B)--(A);
				 \draw[dashed] (C)--(D)--(A);
				\coordinate  (O) at ($(A)!0.5!(C)$);
				 \coordinate (S) at ($(O)+(-0.5,5)$);
				\coordinate  (M) at ($(S)!0.5!(A)$);
				 \draw[dashed] (O)--(S)--(D)--(B)
					 (A)--(C)
					(D)--(M)--(O); 
				\draw (S)--(A) 
					 (S)--(B)--(M)
					(S)--(C); 
				 \tkzMarkRightAngle (A,O,D)
				\tkzMarkRightAngle (S,D,A) 
				\tkzMarkRightAngle  (S,B,A)
				\begin{footnotesize}
					\draw (S)  node[above]{$S$};
					\draw (C)  node[left]{$B$};
					 \draw (D) node[left]{$A$};
					\draw (A)  node[right]{$D$};
					\draw  (B) node[right]{$C$};
					 \draw (M) node[right]{$M$};
					\draw (O) node[below]{$O$}; 
				\end{footnotesize}
			\end{tikzpicture}
		 }
	 }
\end{ex}

% nguồn: chuyen_dang (nhập 22/11/2026)
\begin{bt}%[8H3H10-2]%[ma:7fe5487e]
% nguon: Archive/2024/TOAN_11_XE/TOAN_11_XE/data4/11C2/KNTT-2.OTC.tex #502
 
	Tứ giác $ABCD$ có số đo bốn góc $A,B,C,D$ theo thứ tự lập thành cấp số cộng. Biết số đo góc $C$ gấp $5$ lần số đo góc $A$. Tính số đo các góc của tứ giác $ABCD$ theo đơn vị độ.
	\loigiai{
		Do bốn góc $A,B,C,D$ theo thứ tự lập thành cấp số cộng nên số đo $4$ góc của tứ giác $ABCD$ lần lượt là $u_1,u_1+d,u_1+2d,u_1+3d$. \\
		Theo giả thiết, ta có: $\heva{u_1+u_1+d+u_1+2d+u_1+3d=360 \\ u_1+2d=5u_1}  \Leftrightarrow  \heva{4u_1+6d=360 \\ -4u_1+2d=0} \Leftrightarrow \heva{u_1=22{,}5 \\ d=45.}$\\
		Vậy số đo các góc của tứ giác lần lượt là $22{,}5^\circ; 67{,}5^\circ; 112{,}5^\circ; 157{,}5^\circ$.
	}
\end{bt}

% nguồn: chuyen_dang (nhập 22/11/2026)
\begin{bt}[KNTT, Mức độ 2]%[24-25 Bài giảng K10-K11, Phan Chiến Thắng]%[8H3H10-2]%[ma:eacad5b4]
	Cho tam giác $ABC$ có $\widehat{A}=135^\circ$, $\widehat{C}=15^\circ$ và $b=12$. Tính $a,c,R$ và số đo góc $B$.
	\loigiai{Ta có $\widehat{B}=180^\circ-(\widehat{A}+\widehat{C})=30^\circ$.\\
		Áp dụng định lí sin ta có $\dfrac{a}{\sin 135^\circ}=\dfrac{12}{\sin 30^\circ}=\dfrac{c}{\sin 15^\circ}=2R$. Suy ra $a=\dfrac{12}{\sin 30^\circ}\dot\sin 135^\circ=12\sqrt{2}$.\\
		$c=\dfrac{12}{\sin 30^\circ}\dot\sin 15^\circ\approx 6{,}21$; $R=\dfrac{12}{2\sin 30^\circ}=12$.
	}
\end{bt}

% nguồn: chuyen_dang (nhập 22/11/2026)
\begin{ex}%[8H3V10-2]%[ma:089522db]
	Một tứ giác lồi có số đo các góc lập thành một cấp số nhân. Biết rằng số đo của góc nhỏ nhất bằng $\dfrac{1}{9}$ số đo của góc nhỏ thứ ba. Hãy tính số đo của các góc trong tứ  giác đó.
	 \choice
		{$5^\circ$, $ 15^\circ$, $ 45^\circ$, $  225^\circ$}
		{\True $9^\circ$, $ 27^\circ$, $  81^\circ$, $ 243^\circ$}
		{$7^\circ$, $ 21^\circ$, $ 63^\circ$, $ 269^\circ$} 
		{$8^\circ$, $ 32^\circ$, $  72^\circ$, $ 248^\circ$}
	 \loigiai{
		Gọi số đo các góc của tứ giác là $a$, $aq$, $aq^2$, $aq^3$, trong đó $q >  1$.\\
		Theo giả thiết,  ta có $a=\dfrac{1}{9} aq^2$ nên $q=3$.\\
		Suy ra các góc của tứ giác là  $a$, $3a$, $9a$, $27a$.\\
		Vì tổng các góc trong tứ  giác bằng $360^\circ$ nên ta có $a+3a+9a+27a=360^\circ \Leftrightarrow a=9^\circ$.
	} 
\end{ex}
